\section{Introduction}
\label{sec:intro}


In many data-centric applications --such as those handling large volumes of web or medical information-- explicit reasoning about both structure and content is essential. Classical relational databases often fall short in expressiveness for such use cases, leading to the development of \emph{semi-structured data models}, typically based on labeled trees or graphs~\cite{Bune:semi97,Abit:data99}. A prominent example of such semi-structured data models is XML (eXtensible Markup Language), which encodes structural information via labels from a finite alphabet, and data content using values drawn from an infinite domain. XML documents are commonly queried using XPath, a standard language widely used in tools like XSLT~\cite{xslt:w3c} and XQuery~\cite{xquery:w3c}.  

XPath is, fundamentally, a general purpose language for
addressing, searching, and matching pieces of an XML document. It is an open standard
World Wide Web Consortium (W3C) Recommendation~\cite{xpath:w3c}.
Core-XPath, the fragment of XPath 1.0 that captures its navigational capabilities~\cite{GKP05}, can express structural properties (e.g., tag names, tree traversal), but not conditions involving the actual data values. Logically, Core-XPath corresponds to a classical modal logic~\cite{blackburn2001modal,blackburn06}, and its satisfiability problem is decidable, even under DTD constraints~\cite{M04,BFG08}. Its expressive power has been precisely characterized: it is equivalent to two-variable first-order logic (FO$^2$) on trees~\cite{MdR05} and it is strictly less expressive than PDL with converse~\cite{BK08}. Sound and complete axiomatizations are available in~\cite{CateM09,cateLM10}.

However, Core-XPath’s inability to compare data values across nodes limits its applicability, notably preventing the expression of joins --an essential construct in query languages. To address this, the logic \xpathd (also known as Core-Data-XPath~\cite{BojanczykMSS09}) extends Core-XPath with data-aware features, allowing for expressions such as $\tup{\alpha =_{\textsf{c}} \beta}$ or $\tup{\alpha \neq_{\textsf{c}} \beta}$, where $\alpha$ and $\beta$ are path expressions navigating the XML tree via axes like child, descendant, or sibling. Such expressions enable comparisons of data values for \textsf{c} in nodes reachable via different paths, allowing for significantly richer queries.

By way of example, an expression of the form $\tup{\alpha =_{\textsf{c}} \beta}$ is true at a node $x$ in a data tree if there exist two nodes reachable from $x$ via $\alpha$ and $\beta$ paths respectively, such that the data values for \textsf{c} of $y$ and $z$ are equal. This situation is illustrated in~\Cref{fig:datatree} depicting how the expression {\em ``there is
a one-step descendant via $\textsf{e}$ and a two-steps descendant via $\textsf{e}$ with the same data
value for \textsf{price}''} (written $\tup{\textsf{e} =_{\textsf{price}} \textsf{e}\textsf{e}}$) holds at $x$, given the presence of $u$ and~$z$. The
expression {\em ``there are two two-steps descendant with different data value for \textsf{price}''} (written $\tup{\textsf{e}\textsf{e} \neq_{\textsf{price}} \textsf{e}\textsf{e}}$)
is also true at $x$, because $v$ and $w$ have different data for \textsf{price}.

%\begin{figure}[!h]
%\centering
%\begin{tikzpicture}
%[
%    transform shape, thick,
%    every node/.style = {draw, shape = rectangle, rounded corners, minimum size = 7mm},
%    grow = down,  % alignment of characters
%    level 1/.style = {sibling distance=5cm},
%    level 2/.style = {sibling distance=2.5cm}, 
%    level 3/.style = {sibling distance=3cm}, 
%    level distance = 1.5cm
%  ]
%\node (x) {$a,0$ }
%child {node (y) {$b,2$}  edge from parent[->]
%child {node (u) {$a,1$}}
%child {node (v) {$b,0$}}
%child {node (w) {$a,0$}}}
%child {node (z) {$a,1$}  edge from parent[->]
%};

%\begin{scope}[nodes = {right = 11pt}]
%\node[draw=none]  at (x) {$x$};
%\node[draw=none]  at (y) {$y$};
%\node[draw=none]  at (z) {$z$};
%\node[draw=none]  at (u) {$u$};
%\node[draw=none]  at (v) {$v$};
%\node[draw=none]  at (w) {$w$};
%\end{scope}
%
%\end{tikzpicture}
%\caption{An example of a data tree. Letters are labels, numbers are data.}\label{fig:datatree}
%\end{figure}

\begin{figure}[h!]	
	
		\begin{tikzpicture}[
			>=latex
			]
			
			% Nodes
			\node[draw, shape = rectangle, rounded corners, minimum size = 7mm] 
			   (a) at (0,0) {};
			\node[draw, shape = rectangle, rounded corners, minimum size = 7mm] 
			   (b) [below left=of a] {};
			\node[draw, shape = rectangle, rounded corners, minimum size = 7mm] 
			   (c) [below right=of a] {\textsf{price}=\$10};
			\node[draw, shape = rectangle, rounded corners, minimum size = 7mm] 
			   (d) [below left= of b] {\textsf{price}=\$10};
			\node[draw, shape = rectangle, rounded corners, minimum size = 7mm] 
			   (e) [below=of b] {\textsf{price}=\$10};
			\node[draw, shape = rectangle, rounded corners, minimum size = 7mm] 
			   (f) [below right=of b] {\textsf{price}=\$20};
			
			% Edges with label $e$
			\draw[->] (a) -- (b) node[midway, left] {\textsf{e}};
			\draw[->] (a) -- (c) node[midway, left] {\textsf{e}};
			\draw[->] (b) -- (d) node[midway, above] {\textsf{e}};
			\draw[->] (b) -- (e) node[midway, right] {\textsf{e}};
			\draw[->] (b) -- (f) node[midway, above] {\textsf{e}};
		
		\begin{scope}[nodes = {right = 11pt}]
			\node[draw=none]  at (a) {$x$};
			\node[draw=none]  at (b) {$y$};
		\end{scope}

		\begin{scope}[nodes = {below = 11pt}]
			\node[draw=none]  at (c) {$z$};
			\node[draw=none]  at (d) {$u$};
			\node[draw=none]  at (e) {$v$};
			\node[draw=none]  at (f) {$w$};
		\end{scope}
			
		\end{tikzpicture}
	
	
	\caption{An example of a data tree.}\label{fig:datatree}
\end{figure}


The language of $\xpathd$ allows us to compare data values at the end of a
path, by equality or inequality. However, it does not grant access
to the concrete data value of nodes. This makes it possible to work with an abstraction of data trees:
instead of having concrete data values in each node, we have an
equivalence relation between nodes linking nodes with the same value. In the data tree from~\Cref{fig:datatree}, the relation consists of four equivalence
classes: $\{u,v,z\}$, $\{w\}$, $\{x\}$ and $\{y\}$. This approach treats uniformly nodes with unique data values of a certain type and nodes where data of that type is undefined (see~\cite{arec:data23} for a treatment of incomplete information). 

Recent articles investigate $\xpathd$ from a modal logic perspective.  
For example, the complexity of the satisfiability and evaluation problems of different fragments are discussed
in~\cite{FigPhD,FigueiraS11,Figueira12ACM,Fig13}, while model theory and
expressivity are studied
in~\cite{ADF14,ICDT14,ICDT14Jair,ADF14journal,KR16,GonzalezA21}.  

In this article, we will focus on reasoning methods for $\xpathd$. 
In this line of work, a Gentzen-style sequent calculus is given for a
very restricted fragment of $\xpathd$, named DataGL, in~\cite{BaeldeLS15}. In DataGL, data
comparisons are allowed only between the evaluation point and its
successors. An extension of the equational axiomatic system
from~\cite{cateLM10} is introduced in~\cite{ADFF16report} that is proved sound and complete for \xpathd. This extension allows
downward navigation and equality/inequality tests.
More recently, \cite{ArecF:hilb16,ArecesF21} provide an axiomatization for this logic,
extended with upward navigation, and with nominals and satisfiability
modalities from Hybrid Logic~\cite{arec:hybr05b}. As argued in, e.g., \cite{Brauner07}, features from Hybrid Logic can be effectively leveraged in the design of proof methods for modal logics.  In~\cite{ArecesFS17},
we introduced sound and complete tableau calculus for hybrid $\xpathd$.
%Lastly, an intuitionistic variant of XPath is proposed in~\cite{ArecesCDF23}, aimed at addressing incomplete information and data updates.
%This related work contributes not only to the formal investigation of $\xpathd$, but also sheds new light on a broader family of modal logics featuring data comparisons.


\paragraph{Contributions.}  
In this article we start by introducing a sound, complete and terminating tableau
calculus for $\xpathd$ with downward navigation, where node expressions
are extended with nominals (special labels that are true in only one
node), and path expressions are extended with the satisfiability modality
(allowing the navigation to some particular named node), over the class of all data-graphs. We call this
logic $\Hxpath$. 
%We take inspiration from tools and methods from Hybrid Logic, since proof methods in Hybrid Logic enjoy some nice properties, as stated in~\cite{Brauner07,Brauner2011}. This claim in partly supported by the possibility of internalizing tableaux features into the language. 

Taking inspiration from~\cite{BolanderB07}, we use the hybrid logic operators to provide internalized tableau rules, which we prove are sound and complete for \Hxpath.  The main intuition is
that nominals and satisfiability modalities can be used in tableaux to
keep track of the evaluation of an expression during an attempt to build a
model. This enables us to define tableaux without any help of external features, like labels or relational annotations, as they can be internalized into the language of $\Hxpath$. Crucially, we show how the hybrid machinery can also be used to help with data comparisons. 
We show that the tableau calculus can be ran in polynomial space,
without compromising completeness, establishing in this way that the satisfiability problem
for $\Hxpath$ is $\pspace$-complete.

We then extend our tableaux to handle different structures. First, we show how to enforce certain reachability conditions over the accessibility relations. With this at hand, we provide a method to check satisfiability of $\Hxpath$ over the class of data trees, and show that the complexity of the satisfiability problem remains in \pspace. 
Then, we extend the calculus with so-called pure axioms and node creating rules. These two extensions can be used to characterize many other model classes, and introduce new operators. We show that soundness and  completeness of the resulting tableau calculi follow automatically (even though termination can be compromised). 

This work is an extension of~\cite{ArecesFS17} and improves on it in several ways.  Most importantly we provide details on the termination argument, which was only sketched in the previous article, and show that the $\pspace$ upper bound holds also for trees. To this end, we start by introducing changes to the language that simplify proofs. Concretely, we extend the base language with primitive expressions of the form $\tup{\pathh}\varphi$. This enables us to obtain internalized tableaux, and consequently to prove termination  in a simpler way. Second, we reduce the number of rules by exploring only the left side in a data comparison and adding a commutativity rule. Finally, the current presentation includes a detailed algorithm suitable for implementation.

The contributions in this article serve as a starting point for research and mechanization of other data-aware logics. For example our approach can be use to investigate the logical properties of languages like GQL~\cite{FrancisGGLMMMPR23,FrancisGGLMMMPR23b} and SHACL~\cite{Ortiz23,Ahmetaj0S23}, Modal Logics with Concrete Domains (see, e.g.,~\cite{DemriQ21}), and to study connections with Register Automata (see, e.g.,~\cite{KaminskiF94}).

% We aimed at a very detailed presentation of the different results discussed hereafter, both to make sure that our contributions are self-contained, and also to fill some gaps in previous literature in which, sometimes, too many steps are left to the reader. 

\paragraph{Organization.}  
\Cref{sec:basic} introduces the syntax
and semantics of $\Hxpath$.  \Cref{sec:tableaux} deals with a tableau calculus for
$\Hxpath$, and with examples of how to build tableaux for concrete expressions. This section also shows the calculus is sound and complete with respect to the class of all data models.
\Cref{sec:termination} deals with the proof of termination of the calculus, whereas~\Cref{sec:complexity} shows that the satisfiability problem for $\Hxpath$ is \pspace-complete.
\Cref{sec:extensionplus} extends the calculus to work
over the class of forest and tree models, and discusses completeness and complexity, showing that the complexity of the satisfiability problem over these classes remains \pspace-complete. 
\Cref{sec:pure-ext} introduces extensions of the calculus with pure axioms and node creating rules. These new rules permits the characterization of new classes of models and the definition of new operators. we argue that the resulting tableau calculi remain sound and complete (but termination might be compromised). \Cref{sec:final} includes final remarks and discusses future lines of research.



%%% Local Variables:
%%% mode: latex
%%% TeX-master: "main"
%%% End:
